\chapter{Implementazione software}
\label{cha:implementazione-software}
L'implementazione delle simulazioni è stata sviluppata utilizzando il linguaggio di programmazione Python e la libreria REBOUND, nella versione 5.0.0, per l'integrazione numerica di sistemi gravitazionali a $N$-corpi \cite{rein2012rebound}. REBOUND fornisce gli strumenti necessari per la gestione delle particelle e l'integrazione delle rispettive equazioni del moto. Le funzionalità specifiche del progetto, tra cui la generazione degli ammassi secondo il modello di Plummer, la definizione del potenziale galattico e il calcolo delle grandezze diagnostiche, sono state implementate mediante moduli dedicati, le cui principali funzionalità sono riassunte nella Tabella~\ref{tab:moduli-software}.
\begin{table}[H]
	\centering
	\small
	\caption{Moduli software sviluppati e relative funzionalità.}
	\label{tab:moduli-software}
	\begin{tabular}{p{0.28\textwidth}p{0.66\textwidth}}
		\hline
		\textbf{Modulo} & \textbf{Funzionalità} \\
		\hline
		\texttt{plummer.py} &
		Generazione delle posizioni e delle velocità iniziali e costruzione degli ammassi stellari. \\
		\hline
		\texttt{galaxy.py} &
		Definizione del potenziale galattico, calcolo dell'accelerazione gravitazionale esercitata dalla galassia sulle stelle e determinazione delle condizioni orbitali iniziali. \\
		\hline
		\texttt{simulation.py} &
		Interfaccia con REBOUND e gestione delle stelle, dei parametri di integrazione e dell'evoluzione dinamica. \\
		\hline
		\texttt{cluster\_generation.py} &
		Coordinamento della generazione dell'ammasso e della creazione della relativa simulazione, integrando le funzionalità di \texttt{plummer.py} e \texttt{simulation.py}. \\
		\hline
		\texttt{galactic\_tidal\_stripping.py} &
		Coordinamento della configurazione del sistema galattico e delle condizioni orbitali con la creazione e l'esecuzione della simulazione, integrando le funzionalità di \texttt{galaxy.py} e \texttt{simulation.py}. \\
		\hline
		\texttt{cluster\_diagnostics.py}, \texttt{galactic\_diagnostics.py} &
		Calcolo delle energie, del rapporto viriale, delle stelle legate, della frazione di massa legata e del raggio di metà massa. \\
		\hline
	\end{tabular}
\end{table}
\noindent Al termine delle simulazioni, i risultati sono stati raccolti e organizzati per la costruzione del dataset e per le successive analisi. Per le operazioni numeriche è stata utilizzata NumPy \cite{harris2020array}, per la manipolazione e l'organizzazione dei dati in strutture tabellari pandas \cite{mckinney2010data}, per la produzione dei grafici Matplotlib \cite{hunter2007matplotlib}, mentre l'addestramento e la valutazione dei classificatori SVC e MLP sono stati realizzati mediante scikit-learn \cite{pedregosa2011scikit}.\\
Il codice sorgente, i file di configurazione e le istruzioni per l'esecuzione sono disponibili nel repository pubblico del progetto su GitHub: \url{https://github.com/Qw3rty26/GTS_ML/}